\section{A lower bound for sparse parity at every sample size}
\label{sec:sparse-gaussian}
Let $U\subseteq[d]$, $|U|=k$, and $f=\chi_U$. Write
$F_\theta(S)=\mathbb P_x\{q_S(x)\le\theta\}$ for the actual
Gaussian sign-posterior confidence failure, with
\[
 \theta=\frac{13}{20},\qquad
 c_\theta=\sqrt{2\pi}\left(\frac1e-\frac7{20}\right),\qquad
 a_\theta=\frac{c_\theta}{2\sqrt2}.
\]
The next theorem has a failure fraction independent of both $k$ and $d$.

\begin{theorem}[Ambient-dimension lower bound]
\label{thm:sparse-gaussian}
For each fixed depth $L\ge2$, there are constants $c>0$ and $d_0$
such that, whenever $d\ge d_0$, $1\le k\le d$, and
\[
 0\le n\le\exp\!\left(c k\log\frac{ed}{k}\right),
\]
independent uniform training inputs satisfy
\[
 \mathbb P_S\{F_{13/20}(S)\ge\varepsilon\}\ge\frac34,
 \qquad \varepsilon=\frac{c_\theta^2}{256}>0.
\]
The assertion holds for each specified sample size. In the network model
of the width-limit proposition, every fixed $\gamma<0$ consequently
satisfies the same lower bound for
$\liminf_{N\to\infty}\mathbb P_x\{q_{N,S}(x)\le2/3\}$.
\end{theorem}

\begin{proof}
Fix a failure probability $\delta=1/40$. We choose a sampling constant
$A$ sufficiently large, a degree fraction $\alpha>0$ sufficiently small,
and an odd integer $r_0\ge11$ sufficiently large, in that order.
All these choices depend only on the fixed depth and universal constants.
Further decreases in $\alpha$ are harmless.

We first treat $k\ge K_0$, where the fixed constant $K_0$ is chosen
below; bounded $k$ will be covered at the end.
For a specified $n>0$, let $r$ be the largest odd integer satisfying
\begin{equation}\label{eq:sparse-coarse-cutoff}
 A D_r\log(D_r+1)\le n.
\end{equation}
First suppose this cutoff exists and $r\ge r_0$. By choosing $c$ in the
theorem sufficiently small and then $K_0$ sufficiently large, we ensure
\begin{equation}\label{eq:sparse-cutoff-range}
 r\le\alpha k\le\alpha d,\qquad u:=r+2<k,
 \qquad n<A D_{r+2}\log(D_{r+2}+1).
\end{equation}
Indeed, the binomial lower bound
$\binom dr\ge(d/r)^r$ shows that, for $r\ge\alpha k$, the logarithm
of the sampling threshold is at least a constant multiple of
$\alpha k\log(ed/k)$. Rounding by two can be absorbed by increasing
$K_0$. Also $k\log(ed/k)\le d$, so the same choice of $c$ puts $n$
within the small-exponential range needed for the random Gram bounds.
Uniformly in this shell interval, for large $d$,
\[
 \log(2n)\le C\log(D_r+1)\le C r\log(eD/r),\qquad D=d+1.
\]

Put $\rho=r/D$, $H_\rho=\log(e/\rho)$, $\tau=\tau_r$,
$\ell_r=\lambda_{r+1}$, and
\[
 B=C_{L,\delta}\left(1+
             \frac{\sqrt{nd\,s_{2,r}}}{\tau}
             +\frac{n\ell_r}{\tau}\right),
 \qquad \eta^2=C_0\frac{D_r\log(2n)}n.
\]
Choose $C_0$ as required by Lemma~\ref{lem:multiband-process}.
Lemma~\ref{lem:sparse-gram} and its near-orthogonal coherence estimate
imply, with probability at least $1-\delta$ for all sufficiently large
$d$,
\begin{equation}\label{eq:sparse-low-gram}
 \|Q_{r,SS}\|_{\mathrm{op}}\le B\tau.
\end{equation}
This holds uniformly over $r\ge r_0$ in
\eqref{eq:sparse-cutoff-range}; it does not require $r\gg\log d$.
The sharp spectral ratios give the crucial cancellation
\begin{equation}\label{eq:sparse-variable-accuracy}
 \eta^2B\le C_{L,\delta}\left[
      \frac1A+\sqrt{\frac{\rho H_\rho}{A}}+
                     \rho^2H_\rho\right].
\end{equation}
For completeness, the constant term uses $\eta^2\le C/A$.
The eigenvalue term is bounded by
\[
 \eta^2\frac{n\ell_r}{\tau}
 \le\frac{C D_r\log(2n)}{r\binom D{r+2}}
 \le C\rho^2H_\rho.
\]
The square of the variance term is at most
\[
 \left(\eta^2\frac{\sqrt{nd\,s_{2,r}}}{\tau}\right)^2
 \le\frac CA\frac{C D_r\log(2n)d}{r^2\binom D{r+2}}
 \le\frac CA\rho H_\rho.
\]
Thus $\eta\sqrt B$ can be made arbitrarily small uniformly in $n$
by choosing $A$ large and then $\alpha$ small.

Let $m=\mathbb E[G\mid\mathcal A_S]$ and $p=\Pi_{\le r}m$.
By Lemma~\ref{lem:schedule-mean}, $\|m\|_K^2\le n$ and its training
signs are correct. The ordinary first-power ReLU sampling event in
Lemma~\ref{lem:schedule-relu}, at accuracy $1/4$, has probability at
least $1-\delta$ after increasing $A$. Since $r<k$, Walsh orthogonality
gives $\mathbb E fp=0$, so on this event and
\eqref{eq:sparse-low-gram},
\begin{equation}\label{eq:sparse-head-radius}
 \frac14\mathbb E|p|
 \le\mathbb P_n(-fp)_+
 \le\sqrt{\mathbb P_n(m-p)^2}
 \le\sqrt{B\tau}.
\end{equation}
Set the deterministic radius $R=4\sqrt{B\tau}$.

\medskip\noindent
\emph{Only one additional paired level is required.}
Write
\[
 m=p+b+t,\qquad b=\Pi_{(r,u]}m,\quad t=\Pi_{>u}m,\quad u=r+2,
 \qquad \Delta=\tau_r-\tau_u.
\]
The level-mass estimate and the two-step binomial ratio give
\begin{align}
 \Delta/\tau&\le C_L/r,\label{eq:sparse-band-trace}\\
 n\ell_u/\tau&\le C_L A\rho^2H_\rho,
       \label{eq:sparse-high-eigenvalue}\\
 nd\,s_{2,u}/\tau^2&\le C_L A\rho H_\rho.
       \label{eq:sparse-high-variance}
\end{align}
For example, maximality yields
$n\le C A\binom D{r+2}rH_\rho$, whereas the two high-tail bounds
involve $\binom D{r+4}$ in the denominator; its ratio with
$\binom D{r+2}$ supplies a factor $O(\rho^2)$.
By increasing $r_0$ we have $\tau_u\ge3\tau/4$.
By reducing $\alpha$, Lemma~\ref{lem:sparse-gram} applied at $u$
now gives, with probability at least $1-\delta$,
\begin{equation}\label{eq:sparse-high-gram}
 \|Q_{u,SS}\|_{\mathrm{op}}\le2\tau.
\end{equation}
The coherence remainder tends to zero uniformly, including at fixed
$r_0$. The two displayed high-tail terms can be made as small as needed
before taking $d$ large.

Put $\zeta=n\ell_u/\tau$. We choose $\alpha$ small enough that
\begin{equation}\label{eq:sparse-zeta-choice}
 e_*:=(4+a_\theta^{-2})\zeta\le a_\theta^2/32.
\end{equation}
This is possible uniformly by \eqref{eq:sparse-high-eigenvalue}.
The RKHS bound, the high-tail Gram event, and the variance rank bound
imply
\[
 \mathbb P_n t^2\le2\tau,\qquad
 \|t\|_2^2\le\zeta\tau,
 \qquad
 \mathbb P_x\{\sigma_x^2<\tau/2\}\le4\zeta.
\]
For the last assertion, condition additionally on the full degree-$u$
field. Its remaining diagonal before training is
$\tau_u\ge3\tau/4$, and the normalized trace of the variance reduction
is at most $n\ell_u=\zeta\tau$.
Consequently, except on an input set of measure at most $e_*$, both
$\sigma_x^2\ge\tau/2$ and $|t(x)|\le a_\theta\sqrt\tau$ hold.
The confidence lemma then shows that
\begin{equation}\label{eq:sparse-confidence-margin}
 f(x)(p(x)+b(x))>a_\theta\sqrt\tau
\end{equation}
at every remaining point where $q_S(x)>\theta$.

\medskip\noindent
\emph{A fixed stop-loss bound rules out hiding the negative mass.}
Choose the deterministic threshold
$T=4\sqrt\tau/a_\theta$.
Apply Lemma~\ref{lem:multiband-process} to the deterministic radius
$R=4\sqrt{B\tau}$ and this $T$.
By Markov's inequality its uniform empirical error is at most
\[
 4\delta^{-1}(R\eta+\sqrt\Delta)
 \le a_\theta\sqrt\tau/8
\]
with probability at least $1-\delta$, after the stated choices:
\eqref{eq:sparse-variable-accuracy} controls the first term and
\eqref{eq:sparse-band-trace} controls the second by increasing $r_0$.
The actual $p,b$ belong to this class on the earlier events, since
\eqref{eq:sparse-head-radius} bounds $\mathbb E|p|$ and
$\|b\|_K\le\sqrt n$.
Let $h=p+b$. Training consistency gives $(-f_i h_i)_+\le|t_i|$.
Using $(z-T)_+\le z^2/(4T)$ for $z\ge0$ and
\eqref{eq:sparse-high-gram}, we obtain
\[
 \mathbb P_n(-fh-T)_+
 \le\frac{\mathbb P_n t^2}{4T}
 \le a_\theta\sqrt\tau/8,
 \qquad
 \mathbb E_x(-fh-T)_+\le a_\theta\sqrt\tau/4.
\]
By \eqref{eq:sparse-confidence-margin}, the positive-part integral of
$fh$ is at least
$a_\theta\sqrt\tau(1-F_\theta-e_*)$, and its negative set has measure
at most $F_\theta+e_*$. Crucially, $\deg h\le u<k$, so
$\mathbb E fh=0$ exactly. Equality of the positive and negative
integrals therefore gives
\[
 a_\theta(1-F_\theta-e_*)
 \le\frac4{a_\theta}(F_\theta+e_*)+\frac{a_\theta}4.
\]
It follows that
\[
 F_\theta\ge\frac{3a_\theta^2}{4(4+a_\theta^2)}-e_*
             \ge\frac{a_\theta^2}{32}=\varepsilon,
\]
where we used \eqref{eq:sparse-zeta-choice} and $a_\theta^2<1$.
The low Gram, high Gram, ReLU, and fixed-stop-loss events have combined
probability at least $1-4\delta>3/4$ for all sufficiently large $d$.
No independence between these events is assumed.

\medskip\noindent
\emph{The remaining, smaller samples.}
If the maximal cutoff is below $r_0$, or does not exist, then
$n<A D_{r_0}\log(D_{r_0}+1)$. Fix the odd integer $s=r_0+2$.
The sharp spectrum gives, uniformly in this entire range,
\[
 \frac{n\ell_s}{\tau_s}
 \le C_{L,A,r_0}\frac{D_{r_0}\log(D_{r_0}+1)}
                           {s\binom D{s+2}}=o(1)
 \qquad(d\to\infty).
\]
Let $p=\Pi_{\le s}m$. The mean-norm, residual-variance, and confidence
lemmas show that, except on an $o(1)$ fraction of inputs, confidence
above $\theta$ implies $fp>0$. Conditional on any values of the
$d-k$ nuisance coordinates, $p$ is a degree-$s$ polynomial on the
$k$ target coordinates. Lemma~\ref{lem:ptf-parity} thus bounds its
strictly correct fraction by $1-b_k(s)$, where
\[
 b_k(s)=2^{-k}\sum_{j=0}^{\lfloor(k-s-1)/2\rfloor}\binom kj.
\]
Choose $K_0$ sufficiently large relative to the fixed $s$ that
$b_k(s)\ge2/5$ for every $k\ge K_0$.
This gives $F_\theta\ge2/5-o(1)>\varepsilon$ for every such training
sample. At $n=0$, $q_S=1/2$ identically.
All errors above are uniform over $k\ge K_0$ and the specified sample
range after fixing the constants.

Finally shrink $c$, if necessary, so that $c\le1/(2K_0)$.
For $1\le k<K_0$ the sample bound implies $n\le\sqrt{ed}$.
The degree-zero projection $p_0=\Pi_{\le0}m$ is constant.
The first paired eigenvalues satisfy $\lambda_0=\lambda_1$;
since the total trace is one, $(d+1)\lambda_0\le1$.
Consequently its tail satisfies the exact bound
$\ell_0/\tau_0=\lambda_1/(1-\lambda_0)\le1/d$.
The same mean and residual-variance estimates show that, outside a
fraction $O_\theta(n/d)=o(1)$ of inputs, confidence above $\theta$
implies $fp_0>0$. A constant can be strictly correct on at most half
the cube for any nonconstant parity. Therefore
$F_\theta\ge1/2-o(1)>\varepsilon$ for every training sample in this
remaining range. A single sufficiently large $d_0$ covers all cases.

Finally, the Gaussian sign-posterior width limit holds at fixed
$d,k,S$ for every fixed $\gamma<0$. At any point with
$q_S(x)\le13/20$, finite-width correct-label probabilities are eventually
below $2/3$. Since the cube is finite, the asserted $\liminf$ transfer
holds simultaneously at all these points.
\end{proof}
